\documentclass[12pt]{article}

\usepackage[english]{babel}
\usepackage[letterpaper,top=2.5cm,bottom=2.5cm,left=3cm,right=3cm]{geometry}
\usepackage{amsmath}
\usepackage{amssymb}
\usepackage{graphicx}
\usepackage{float}
\usepackage{booktabs}
\usepackage[colorlinks=true, allcolors=blue]{hyperref}
\usepackage{setspace}
\doublespacing

\DeclareMathOperator{\sgn}{sgn}
\newcommand{\RE}{\mathrm{Re}}
\newcommand{\IM}{\mathrm{Im}}
\newcommand{\infint}{\displaystyle\int_{-\infty}^{+\infty}}

\newcommand{\citein}[1]{\textit{#1}}

% ====================================================================
\title{Supplementary Appendices\\
\large An analytical cyclo-steady-state model for inter-rifting deformation at magmatic rifts}

\date{}

% ====================================================================
\begin{document}
\maketitle

\appendix

% ====================================================================
\section{Fourier reduction and the characteristic velocity}
\label{app:fourier}

Equation~(1) (the governing PDE) is non-local in $x$ and not tractable in physical
space. A spatial Fourier transform resolves this, because the Stokes convolution
diagonalizes in wavenumber space. Defining the transform pair
\begin{equation*}
\hat{u}(k,t) = \int_{-\infty}^{+\infty}u(x,t)e^{-ikx}\,dx,
\qquad
u(x,t) = \frac{1}{2\pi}\int_{-\infty}^{+\infty}\hat{u}(k,t)e^{ikx}\,dk,
\end{equation*}
the transform of $\partial^2 u/\partial x^2$ is $-k^2\hat{u}$. The transform
of the Stokes kernel convolution is the classical result
(Savage and Prescott, 1978):
\begin{equation}
\mathcal{F}\!\left\{
  \int_{-\infty}^{+\infty}\frac{2\eta}{\pi(x-x')^2}
  \frac{\partial u}{\partial t}\,dx'
\right\} = -2\eta|k|\,\frac{\partial\hat{u}}{\partial t}.
\label{eq:stokes_FT_app}
\end{equation}
Applying the Fourier transform to the governing PDE and using
(\ref{eq:stokes_FT_app}):
\begin{equation}
-\frac{EH}{1-\nu^2}k^2\hat{u} - 2\eta|k|\frac{\partial\hat{u}}{\partial t} = 0.
\label{eq:fourier_pre_app}
\end{equation}
Rearranging and introducing the characteristic velocity
\begin{equation}
a := \frac{EH}{2\eta(1-\nu^2)},
\label{eq:a_app}
\end{equation}
equation (\ref{eq:fourier_pre_app}) simplifies to a decoupled first-order ODE for
each wavenumber $k$:
\begin{equation}
\frac{\partial\hat{u}}{\partial t} = -a\,|k|\,\hat{u}.
\label{eq:fourier_app}
\end{equation}
The parameter $a$ encodes the full elastic--viscous coupling of the system:
large $a$ (stiff crust, low viscosity) implies rapid redistribution of
post-diking strain over broad spatial scales; small $a$ implies slow, locally
confined relaxation. In physical space, the operator $|k|$ corresponds to the
fractional half-Laplacian $(-\Delta)^{1/2}$, so (\ref{eq:fourier_app}) is
equivalently written as the fractional diffusion equation
(Metzler and Klafter, 2000; Samko and Ross, 1993; Samko et al., 1993)
\begin{equation}
\frac{\partial u}{\partial t} = -a\,(-\Delta)^{1/2}u.
\end{equation}
This contrasts with classical Fickian diffusion: the half-order operator, a
direct consequence of the non-local Stokes kernel, produces Lorentzian rather
than Gaussian spatial profiles, with tails that decay as $1/x$ rather than
exponentially. The solution of (\ref{eq:fourier_app}) for arbitrary initial data is:
\begin{equation}
\hat{u}(k,t) = \hat{u}(k,0)\,e^{-a|k|t}.
\label{eq:fourier_sol_app}
\end{equation}

% ====================================================================
\section{Step response: Sokhotski--Plemelj and Dirichlet evaluation}
\label{app:step}

A symmetric opening of $2\omega$ at $t = 0$ displaces the right side of the rift
by $+\omega$ and the left by $-\omega$. Exploiting the antisymmetry, the initial
condition is a Heaviside displacement:
\begin{equation}
u(x,0^+) = \omega\,\mathcal{H}(-x).
\label{eq:IC_app}
\end{equation}
Applying the Fourier transform to (\ref{eq:IC_app}) gives
$\hat{u}(k,0) = \omega\int_{-\infty}^{0}e^{-ikx}\,dx$,
which for real $k$ is treated in the distributional sense by adding a
convergence regulator $\varepsilon > 0$:
\begin{equation*}
\int_{-\infty}^{0}e^{-ikx}\,dx
= \lim_{\varepsilon\to 0}\int_{-\infty}^{0}e^{(-ik+\varepsilon)x}\,dx
= \lim_{\varepsilon\to 0}\frac{1}{\varepsilon - ik}.
\end{equation*}
The Sokhotski--Plemelj theorem gives $\pi\delta(k) + i/k$ in the limit, so
the initial Fourier spectrum is:
\begin{equation}
\hat{u}(k,0) = \omega\pi\delta(k) + \frac{i\omega}{k}.
\label{eq:IC_FT_app}
\end{equation}
Substituting (\ref{eq:IC_FT_app}) into (\ref{eq:fourier_sol_app}) and noting that
$\delta(k)$ is active only where $e^{-a|k|t}=1$:
\begin{equation}
\hat{u}(k,t) = \omega\pi\delta(k) + \frac{i\omega}{k}\,e^{-a|k|t}.
\label{eq:evolved_app}
\end{equation}
The inverse Fourier transform splits into two contributions,
$u(x,t) = I_1 + I_2$. The delta-function term gives a constant background
$I_1 = \omega/2$; for $I_2$, the parity of $\sin(kx)/k$ yields:
\begin{equation}
I_2 = -\frac{\omega}{\pi}\int_0^{\infty}\frac{\sin(kx)}{k}\,e^{-akt}\,dk.
\end{equation}
This is a standard Dirichlet-type integral. Applying the complementary
arctangent identity $\arctan(z)+\arctan(1/z)=\pi/2$, the step displacement is:
\begin{equation}
u^{\text{STEP}}(x,t) = \frac{\omega}{\pi}\arctan\!\left(\frac{at}{x}\right).
\end{equation}
Differentiating with respect to $t$, the step velocity is a Lorentzian
profile that broadens in time:
\begin{equation}
v^{\text{STEP}}(x,t) = \frac{\omega a}{\pi}\,\frac{x}{x^2 + a^2t^2}.
\end{equation}
Both expressions satisfy the correct limits: $u^{\text{STEP}}\to\omega/2$ and
$v^{\text{STEP}}\to 0$ as $t\to\infty$, consistent with far-field plate
divergence. The velocity peak migrates outward as $at$, with amplitude decaying
as $1/t$.

% ====================================================================
\section{CSS Digamma representation}
\label{app:digamma}

The CSS velocity arises by extending the finite sum over past intrusions to
an infinite history. Introducing $\tau\in[0,T)$ as the time elapsed since the
most recent intrusion:
\begin{equation}
v_{\text{CSS}}(x,\tau) = \frac{\omega a}{\pi}
  \sum_{j=1}^{\infty}\frac{x}{x^2 + a^2(\tau-jT)^2}.
\label{eq:css_sum_app}
\end{equation}
Because this series converges absolutely, it admits a closed-form expression
via the Digamma function $\psi(z)$. Using the partial-fraction identity
\begin{equation*}
\frac{x}{x^2+a^2(\tau-jT)^2}
= \frac{1}{2i}\!\left[
    \frac{1}{a(\tau-jT)-ix} - \frac{1}{a(\tau-jT)+ix}
  \right],
\end{equation*}
factoring $-aT$ from each denominator and substituting into
(\ref{eq:css_sum_app}), the double sum takes the form of differences of Digamma
series. Using the representation $\psi(z+1) = -\gamma +
\sum_{n=1}^\infty(1/n - 1/(n+z))$, the Euler--Mascheroni constant $\gamma$ and
the $1/n$ terms cancel when the two conjugate series are differenced. With
$Z = \tau/T + ix/(aT)$ and the identity
$\psi(Z)-\psi(\bar{Z}) = 2i\,\IM\{\psi(Z)\}$,
the CSS velocity collapses to:
\begin{equation}
v_{\text{CSS}}(x,\tau) =
  \frac{2\omega}{\pi T}\,\IM\!\left\{
    \psi\!\left(\frac{\tau}{T}+i\frac{x}{aT}\right)
  \right\}.
\end{equation}
Integrating over $\tau$, and using the antiderivative
$d[\ln\Gamma(z)]/dz = \psi(z)$, the CSS displacement is:
\begin{equation}
u_{\text{CSS}}(x,\tau) =
  \frac{2\omega}{\pi}\,\IM\!\left\{
    \ln\Gamma\!\left(\frac{\tau}{T}+i\frac{x}{aT}\right)
    -\ln\Gamma\!\left(i\frac{x}{aT}\right)
  \right\}.
\end{equation}
The far-field conservation property $\int_0^T v_{\text{CSS}}\,d\tau = \omega$
follows from $\Gamma(1+iq) = iq\,\Gamma(iq)$, which implies
$\IM\{\ln\Gamma(1+iq)-\ln\Gamma(iq)\} = \pi/2$, giving the integral $\omega$.

% ====================================================================
\section{Maxwell visco-elastic kernel and governing equation}
\label{app:maxwell}

\subsection*{Elastic Stokes kernel and the correspondence principle}

For a purely elastic half-space the basal traction relates to surface
displacement through a kernel proportional to $\mu_m$:
\begin{equation}
\sigma_{xz}^{\text{BOT, el}}(x) =
  \int_{-\infty}^{+\infty}\frac{2\mu_m}{\pi(x-x')^2}\,u(x')\,dx'.
\label{eq:elastic_kernel_app}
\end{equation}
The Maxwell constitutive relation for deviatoric stress $\tau_{ij}$ and
strain $\varepsilon$ is:
\begin{equation}
\dot{\varepsilon} = \frac{\dot{\tau}}{2\mu_m} + \frac{\tau}{2\eta}.
\end{equation}
To adapt (\ref{eq:elastic_kernel_app}) we employ the elastic--visco-elastic
correspondence principle (Hetland and Hager, 2005; Barbot and Fialko, 2010):
in the Laplace domain, static moduli are replaced by their complex,
frequency-dependent counterparts. Applying the Laplace transform
$\tilde{f}(s)=\mathcal{L}\{f(t)\}$ with zero initial conditions gives:
\begin{equation}
s\tilde{\varepsilon} = \tilde{\tau}\!\left(\frac{s}{2\mu_m}+\frac{1}{2\eta}\right).
\end{equation}
Solving for stress gives $\tilde{\tau} = 2\tilde{\mu}(s)\tilde{\varepsilon}$,
where the effective frequency-dependent shear modulus of the Maxwell body is:
\begin{equation}
\tilde{\mu}(s) = \frac{\mu_m\eta\,s}{\mu_m + \eta s}
               = \mu_m\frac{\tau_R s}{1+\tau_R s},
\end{equation}
with the Maxwell relaxation time $\tau_R := \eta/\mu_m$.

\subsection*{Visco-elastic Stokes kernel in time domain}

Applying the correspondence principle (replacing $\mu_m$ with $\tilde{\mu}(s)$
in the Laplace domain) and factoring out $\eta = \mu_m\tau_R$:
\begin{equation}
(1+\tau_R s)\,\tilde{\sigma}_{xz}^{\text{BOT}}(x,s) =
  \int_{-\infty}^{+\infty}\frac{2\eta}{\pi(x-x')^2}
  \left[s\tilde{u}(x',s)\right]dx'.
\end{equation}
Inverting to the time domain (multiplication by $s$ in the Laplace domain
corresponds to $\partial/\partial t$) yields the visco-elastic Stokes kernel:
\begin{equation}
\sigma_{xz}^{\text{BOT}} + \tau_R\frac{\partial\sigma_{xz}^{\text{BOT}}}{\partial t}
= \int_{-\infty}^{+\infty}\frac{2\eta}{\pi(x-x')^2}
  \frac{\partial u(x',t)}{\partial t}\,dx'.
\label{eq:ve_kernel_app}
\end{equation}
In the limit $\tau_R\to 0$, (\ref{eq:ve_kernel_app}) reduces to the purely viscous
Stokes kernel.

\subsection*{Governing equation in Fourier space}

The equilibrium equation expresses the basal traction in terms
of the elastic crust displacement:
\begin{equation}
\sigma_{xz}^{\text{BOT}} = -\frac{EH}{1-\nu^2}\frac{\partial^2 u}{\partial x^2}.
\end{equation}
Substituting this expression and its time derivative into
(\ref{eq:ve_kernel_app}), then applying the spatial Fourier transform
using (\ref{eq:stokes_FT_app}):
\begin{equation}
\frac{EH}{1-\nu^2}k^2\hat{u}
+ \tau_R\frac{EH}{1-\nu^2}k^2\frac{\partial\hat{u}}{\partial t}
= -2\eta|k|\frac{\partial\hat{u}}{\partial t}.
\end{equation}
Dividing by $2\eta$, using (\ref{eq:a_app}), and collecting
$\partial\hat{u}/\partial t$ terms gives the visco-elastic analogue of
(\ref{eq:fourier_app}):
\begin{equation}
\frac{\partial\hat{u}}{\partial t}
= -\left(\frac{a|k|}{1+\tau_R a|k|}\right)\hat{u}.
\label{eq:maxwell_fourier_app}
\end{equation}
The denominator $1+\tau_R a|k|$ acts as a low-pass filter on the relaxation
rate, slowing the decay of high-wavenumber (near-field) features relative to
the purely viscous case.

\subsection*{Initial condition evolution and step response}

Equation (\ref{eq:maxwell_fourier_app}) is separable in $k$, yielding a decaying
exponential for each wavenumber:
\begin{equation}
\hat{u}(k,t) = \hat{u}(k,0)\exp\!\left(-\frac{a|k|}{1+\tau_R a|k|}\,t\right).
\end{equation}
Applying the Heaviside initial condition (\ref{eq:IC_FT_app}):
\begin{equation}
\hat{u}(k,t) = \omega\pi\delta(k)
  + \frac{i\omega}{k}\exp\!\left(-\frac{a|k|}{1+\tau_R a|k|}\,t\right).
\end{equation}
Taking the inverse Fourier transform and following the same parity argument as
in Appendix~\ref{app:step}, the Maxwell step displacement is:
\begin{equation}
u^{\text{MAX}}(x,t) = \frac{\omega}{2}
  - \frac{\omega}{\pi}\int_0^{+\infty}
    \frac{\sin(kx)}{k}\exp\!\left(-\frac{akt}{1+\tau_R ak}\right)dk,
\end{equation}
since the rational exponent prevents the closed-form Dirichlet evaluation of
Appendix~\ref{app:step}. The surface velocity follows by differentiating with
respect to $t$:
\begin{equation}
v^{\text{MAX}}(x,t) = \frac{\omega a}{\pi}\int_0^{+\infty}
  \frac{\sin(kx)}{1+\tau_R ak}
  \exp\!\left(-\frac{akt}{1+\tau_R ak}\right)dk.
\end{equation}
The exponent reveals two physically distinct regimes. In the long-wavelength
limit ($k\to 0$ or $t\to\infty$), $1+\tau_R ak \approx 1$ and the exponential
reduces to $\exp(-akt)$, recovering the purely viscous Dirichlet integral of
Appendix~\ref{app:step}. In the short-wavelength limit ($k\to\infty$ or
$t\to 0$), $1+\tau_R ak \approx \tau_R ak$ and the exponent saturates:
\begin{equation}
\lim_{k\to\infty}\left(-\frac{akt}{1+\tau_R ak}\right) = -\frac{t}{\tau_R}.
\end{equation}
At early times $t\ll\tau_R$, the exponential approaches unity and no near-field
deformation occurs outside the rift interface. For $t\sim\tau_R$, near-field
deformation decays exponentially at the fixed rate $1/\tau_R$, representing
pure Maxwell relaxation independent of the elastic crust.

% ====================================================================
\section{Maxwell Cyclo-Steady State: geometric series and periodicity condition}
\label{app:maxwell_css}

The Maxwell CSS is the superposition of step-response velocities from all
preceding intrusions:
\begin{equation}
  v_{\text{CSS}}^{\text{MAX}}(x,\tau)
  = \sum_{j=0}^{\infty} v^{\text{MAX}}(x,\,\tau + jT).
\end{equation}
Writing $v^{\text{MAX}}$ in Fourier space as
$\frac{2\omega a}{\pi}\int_0^\infty \sin(kx)\,e^{-\alpha(k)(\tau+jT)}\,dk$,
with the effective Maxwell decay rate
\begin{equation}
  \alpha(k) := \frac{ak}{1+\tau_R ak},
  \label{eq:alpha_k_app}
\end{equation}
each past intrusion contributes a factor $e^{-\alpha(k)jT}$. The series
converges for all $k>0$ and the geometric sum
$\sum_{j=0}^{\infty}e^{-\alpha(k)jT} = (1-e^{-\alpha(k)T})^{-1}$ can be
performed inside the Fourier integral. The same result follows directly
from the requirement that each Fourier mode satisfy the periodicity
condition $\hat{u}(k,\tau+T) = \hat{u}(k,\tau) + i\omega/k$, which
admits the unique solution
\begin{equation}
  \hat{u}(k,\tau) = \frac{i\omega/k}{1-e^{-\alpha(k)T}}\,e^{-\alpha(k)\tau},
  \qquad \tau\in[0,T).
\end{equation}
Inverting gives the Maxwell CSS velocity and displacement:
\begin{align}
  v_{\text{CSS}}^{\text{MAX}}(x,\tau)
  &= \frac{2\omega a}{\pi}\int_0^{\infty}
     \frac{\sin(kx)}{1+\tau_R ak}\,
     \frac{e^{-\alpha(k)\tau}}{1-e^{-\alpha(k)T}}\,dk, \\
  u_{\text{CSS}}^{\text{MAX}}(x,\tau)
  &= \frac{2\omega}{\pi}\int_0^{\infty}
     \frac{\sin(kx)}{k}\,
     \frac{1-e^{-\alpha(k)\tau}}{1-e^{-\alpha(k)T}}\,dk.
\end{align}
Integrating the velocity over one full period, using
$\int_0^T e^{-\alpha\tau}/(1-e^{-\alpha T})\,d\tau = 1/\alpha$ and the
algebraic identity $a/[(1+\tau_R ak)\,\alpha(k)] = 1/k$, gives
\begin{equation}
  \int_0^T v_{\text{CSS}}^{\text{MAX}}(x,\tau)\,d\tau
  = \frac{2\omega}{\pi}\int_0^{\infty}
    \frac{\sin(kx)}{k}\,dk = \omega,
\end{equation}
where the last step is the Dirichlet integral. The full rift half-opening
is recovered every cycle, independent of $\tau_R$. In the limit
$\tau_R\to 0$, $\alpha(k)\to ak$ and $(1+\tau_R ak)^{-1}\to 1$, and the
Maxwell CSS velocity reduces to the viscous CSS Fourier integral whose
closed form is the Digamma series of Appendix~\ref{app:digamma}.

% ====================================================================
\end{document}
